\section{1907? Hjelmslev i odwzorowanie na koło} % lang-pl
\label{sekcja_hjelmslev}
Model w kole da się otrzymać także na inny sposób: bez rzutowań, samym odwzorowaniem płaszczyzny Łobaczewskiego. % lang-pl
W tej sekcji poznamy duńskiego geometrę, którego twierdzenie o środkach odcinków to odwzorowanie umożliwi, i zobaczymy, jak ta sama konstrukcja działa w geometrii euklidesowej. % lang-pl

Johannes Trolle Hjelmslev urodzi się jako Johannes Trolle Petersen, a nazwisko Hjelmslev przyjmie, żeby odróżnić się od Juliusa Petersena.

\begin{theorem}[Hjelmsleva] % lang-pl
\index{twierdzenie!Hjelmsleva} % lang-pl
	Środki odcinków łączących odpowiadające sobie punkty dwóch przystających ciągów punktów (na płaszczyźnie euklidesowej lub Łobaczewskiego) są współliniowe albo pokrywają się. % lang-pl
\end{theorem}

Twierdzenie to znane będzie od dawna, ale dopiero w 1907 roku Hjelmslev odkryje, że należy ono do geometrii absolutnej i że ma liczne zastosowania w geometrii Łobaczewskiego; dowód: Eves \cite[s. 311--312]{eves1_1972}. % lang-pl

Ustalmy punkt $O$ i kąt ostry $\theta$. % lang-pl
Przekształcenie $T$ przypisuje punktowi $P$ taki punkt $P'$, że $\angle POP' = \theta$ i $\angle OP'P = \tfrac{\pi}{2}$, a następnie obraca $P'$ wokół $O$ o $-\theta$ \cite[s. 312]{eves1_1972}. % lang-pl

\begin{proposition}
	Na płaszczyźnie euklidesowej $T$ jest jednokładnością o środku $O$ i skali $\cos\theta$, a więc przekształca płaszczyznę na siebie. % lang-pl
	Na płaszczyźnie Łobaczewskiego $T$ przekształca całą płaszczyznę na wnętrze koła o środku $O$ i promieniu $h$, gdzie $\Pi(h) = \theta$. % lang-pl
\end{proposition}

Dowody: Eves \cite[s. 312--313]{eves1_1972}. % lang-pl
Eves zauważa też, że obraz jest podobny do modelu z sekcji \ref{sekcja_punkty_idealne}, i nie rozwija tego związku z braku miejsca \cite[s. 313]{eves1_1972}. % lang-pl

Hjelmslev będzie też jednym z autorów dowodów twierdzenia o jednorodności (zobacz sekcję \ref{sekcja_suma_katow}): według przypisu Hilberta, cytowanego przez Greenberga, dowód pierwszy poda Dehn, a późniejsze F.~Schur i Hjelmslev \cite[s. 208]{greenberg_2010}. % lang-pl

Zadania: Eves \cite[s. 314]{eves1_1972}, w tym wyprowadzenie z odwzorowania $T$ przenoszalności, symetrii i przechodniości równoległości. % lang-pl
% Źródła: Eves s. 314--317 (§7.9); Wikipedia: Henri Poincaré, János Bolyai.

















































Tekst sekcji geometrie nieeuklidesowe. % lang-pl

